% =====================================================================
%  TESLİM PAKETİ v2 — RENEWABLE ENERGY · TÜRKÇE PARALEL
%
%  ESAS SÜRÜM İNGİLİZCE: manuscript/en/RE_manuscript_v2.tex
%  Bu dosya Gökhan'ın takibi ve ortak çalışma içindir; içerik farkı
%  olursa İngilizce geçerlidir. Dergiye İngilizce sürüm gider.
%
%  DURUM ETİKETLERİ
%    [YENİ]    bu sürümde tam yazıldı
%    [TAŞI]    metin taslak v2'de (.docx) var; oradan aktarılacak.
%              Tablolardaki sayılar .docx'ten birebir çıkarıldı,
%              elle yeniden yazılmadı.
%    [HESAPLA] henüz yapılmamış bir koşu gerektiriyor
% =====================================================================

\documentclass[11pt,a4paper]{article}
\usepackage[a4paper,margin=2.5cm]{geometry}
\usepackage{fontspec}
\usepackage{polyglossia}
\setmainlanguage{turkish}
\setmainfont{TeX Gyre Termes}
\usepackage{amsmath,amssymb}
\usepackage{booktabs,array,longtable,graphicx}
\usepackage[hidelinks]{hyperref}
\usepackage{url}
\setlength{\parindent}{0pt}
\setlength{\parskip}{5pt}

\newcommand{\hesapla}[1]{\textbf{[HESAPLA: #1]}}
\newcommand{\tasi}[1]{\textsf{\small [TAŞI -- #1]}}

\begin{document}

\begin{center}
{\bfseries\large Türbin düzeyinde sızıntısız kısa vadeli rüzgâr gücü tahmini:
düşük gecikmeli nedensel filtre bankası, eğitim gerektirmeyen bir nedensellik
sertifikası ve dürüst değerlendirmenin bedeli\par}
\vspace{6pt}
Gökhan Gülmez\textsuperscript{1}, Kadir Sarıkaya\textsuperscript{2}\par
{\small \textsuperscript{1}Bilgisayar Mühendisliği Bölümü,
\textsuperscript{2}Endüstri Mühendisliği Bölümü,\\
Mühendislik ve Mimarlık Fakültesi, Tokat Gaziosmanpaşa Üniversitesi, Tokat\par}
\vspace{4pt}
{\footnotesize\textsf{Hedef dergi: Renewable Energy (Elsevier, elsarticle).
Yazar sırası gönderimden önce netleştirilecek.}}
\end{center}

% =====================================================================
\section*{ÖNE ÇIKANLAR (Highlights)}
% Elsevier kuralı: 3--5 madde, her biri boşluklar dahil en çok 85 karakter.
% =====================================================================
\begin{itemize}\itemsep1pt
\item Tam seri VMD, altı SDWPF türbininde nedensel MAE'yi \%34--57 eksik gösteriyor
\item Eğitim gerektirmeyen sertifika, nedensel ve sızdıran operatörleri ayırıyor
\item Kilitli nedensel filtre bankası, kayan VMD'yi 1300 kat düşük gecikmeyle yakalıyor
\item Nedensel protokolde ayrıştırma, ham SCADA girdilerine hiçbir şey eklemiyor
\item Oracle-normalize oran, kabul edilen sızıntının ne kadarının gerçekleştiğini veriyor
\end{itemize}

% =====================================================================
\section*{ÖZET}
% [YENİ] — 248 kelimelik İngilizce özetin karşılığı. Renewable Energy
% için operasyonel sonuç (konuşlandırılabilir doğruluk, güncelleme
% maliyeti) öne alındı.
% =====================================================================

Ayrıştırma tabanlı hibritler kısa vadeli rüzgâr gücü tahmini literatürüne hâkim
durumdadır, ancak çoğunlukla güç serisinin tamamı eğitim/test bölmesinden önce
ayrıştırılarak değerlendirilmektedir; bu durumda bir test anını betimleyen
öznitelikler o anın kendi geleceğini kodlar. Dürüst değerlendirme, ayrıştırmanın
her tahmin anında yalnızca geçmiş veriden yeniden hesaplanmasını gerektirir ve
bu hesaplama pahalıdır --- bu yüzden nadiren yapılır. Bu çalışma söz konusu
yapaylığı nicelleştirmekte ve türbin düzeyi SCADA tahmini için maliyet engelini
kaldırmaktadır. SDWPF çiftliğinin altı türbininde, 0,5--6 saatlik ufuklarda ve
sekiz rastgele başlangıçta, tam seri varyasyonel mod ayrıştırması (VMD) aynı
nedensel boru hattının ortalama mutlak hatasını \%34--57 eksik göstermektedir;
fark Holm düzeltmeli Diebold--Mariano testleriyle her türbinde anlamlıdır ve
tahminci mutlak hata kaybıyla eğitildiğinde de değişmemektedir. Bir ileri-bakış
doz--tepki deneyi farkı genlik, faz veya sınır etkilerine değil gelecek
bilgisine atfetmektedir. İki araç tanıtıyoruz: operatörü gelecek bozmaları ve
önek kesmesiyle yoklayan, geçti/kaldı hükmü yerine bir sızıntı erişimi ve
kazancı raporlayan, eğitim gerektirmeyen bir nedensellik sertifikası; ve bir
operatörün kabul ettiği sızıntının ne kadarını fiilen gerçekleştirdiğini ifade
eden oracle-normalize bir oran. Ardından, bantları eğitim döneminden kilitlenen
nedensel bir spektral filtre bankasının, kayan VMD'den istatistiksel olarak
ayırt edilemeyen nedensel doğruluğa yaklaşık 1300 kat düşük güncelleme
gecikmesiyle ulaştığını ve böylece sızıntısız değerlendirmeyi çiftlik ölçeğinde
uygulanabilir kıldığını gösteriyoruz. Nedensel protokol altında ayrıştırma, ham
SCADA girdilerine göre tutarlı bir doğruluk kazancı sağlamamaktadır.

\textbf{Anahtar kelimeler:} rüzgâr gücü tahmini; SCADA; zamansal veri sızıntısı;
varyasyonel mod ayrıştırması; nedensel ön işleme; değerlendirme protokolü;
Diebold--Mariano testi.

% =====================================================================
\section{Giriş}
% =====================================================================
% [YENİ]

\subsection{Operasyonel motivasyon}

Kısa vadeli rüzgâr gücü tahminleri birim taahhüdünü, yedek tedarikini, gün içi
ticareti ve kısıtlama kararlarını yönlendirir; dolayısıyla bir tahmincinin
raporladığı doğruluk değeri akademik bir nicelik değildir, işletmecinin yedek
boyutlandırdığı temeldir. Konuşlandırmada yeniden üretilemeyen bir değer yedeği
yanlış boyutlandırır ve hata tek yönlüdür: olduğundan iyi görünen doğruluk,
olduğundan küçük yedeğe yol açar.

Yayımlanmış kısa vadeli rüzgâr gücü tahmincileri arasında ayrıştırma tabanlı
hibritler en yaygın tasarımdır: gözlenen güç serisi frekans bileşenlerine
ayrılır ve öğrenici bu bileşenler üzerinde eğitilir. Qian ve ark.~[1] bu aileyi
derlemiştir. Tasarım, nadiren belirtilen bir değerlendirme seçimine dayanır:
ayrıştırma \emph{ne zaman} uydurulmaktadır? VMD~[2] ve ampirik mod ailesi
küresel, nedensel olmayan dönüşümlerdir. Biri test dönemini içeren bir seriye
bir kez uydurulursa, $t$ anındaki bileşen değeri $t$ sonrasındaki gözlemlerin
fonksiyonu olur. Hedef modele hiç gösterilmez, ama hedef hakkındaki bilgi
gösterilir.

\subsection{Dürüst protokol neden benimsenmiyor}

Çözüm bilinmektedir: ayrıştırmayı her tahmin anında yalnızca geçmiş veriden
yeniden hesaplamak. Chen ve ark.~[3] sorunu ``sınır sorunu'' olarak kataloglar
ve çözümleri algoritma, örnekleme ve yineleme tabanlı diye sınıflar. Engel
maliyettir. Kayan yeniden-ayrıştırma her tahmin anında yeni bir varyasyonel
optimizasyon gerektirir; Jeon ve ark.~[4] bunun ``gerçek bir hesaplama maliyeti''
taşıdığını ve bu maliyetin ``konuşlandırmayı yansıtan bir değerlendirmenin
bedeli'' olduğunu açıkça yazmaktadır. SDWPF büyüklüğünde bir çiftlikte
(134 türbin) bu bedel karşılanamaz ve literatür bedeli ödememe yoluna gitmiştir.

Bu makale maliyeti bir çekince olarak değil bir mühendislik problemi olarak ele
alır. Temel pratik iddiası, bedelin nedensel doğrulukta ölçülebilir bir kayıp
olmaksızın yaklaşık üç mertebe düşürülebileceği ve böylece dürüst protokolün
çiftlik ölçeğinde kullanılabilir hâle geldiğidir.

\subsection{Bu çalışmanın konumu}

Yapaylığı ilk gözlemleyen ya da ilk ölçen olduğumuzu iddia etmiyoruz ve okurun
neyin zaten yerleşik olduğunu bilmesi gerekir. Kayan ya da örnek bazlı
ayrıştırma benzin tüketimi~[5], ham petrol~[6], elektrik yükü~[7,8], rüzgâr
hızı~[9--13] ve çok kısa vadeli rüzgâr gücü~[14] için önerilmiştir.
Nicelleştirilmiş karşılaştırmalar da vardır: Xu ve Wang~[6] aynı tahmincilerle
kayan ve tüm seri EEMD'yi karşılaştırır; Zhang ve ark.~[15] altı veri kümesi ve
üç ayrıştırmada anormal hata düşüşünün ağırlıkla yüksek frekanslı bileşenlerce
taşındığını raporlar; Zhou ve ark.~[13] rüzgâr hızı tahmininde şişmeyi mutlak
hatada \%40--55, kare hatada \%65--85 olarak ölçer; Jeon ve ark.~[4] Alman yük,
güneş ve rüzgâr üretimi ile bir kıyas serisinde tam seri protokolün aynı
nedensel boru hattının hatasını \%30--93 eksik gösterdiğini bildirir ve
yapaylığı geçmiş uzunluğuna değil ileri-bakışa atfeder.
Valdivieso ve ark.~[16] yük için benzer bir büyüklük raporlar.

Bu külliyatın açık bıraktığı üç şey vardır ve bu makalenin katkısı odur.

\emph{Birincisi, çözümün maliyeti adlandırılmış ama ele alınmamıştır.} Yukarıdaki
çalışmaların hiçbiri nedensel ayrıştırmanın güncelleme başına gecikmesini ölçmez
ya da düşürmez. Bölüm~4.1 kilitli nedensel spektral filtre bankasını tanımlar,
Bölüm~6.5 gecikmesini kayan VMD'ye karşı aynı seride ölçer.

\emph{İkincisi, bir operatörün nedenselliği için uygulanabilir bir test yoktur.}
Jeon ve ark.~[4] beş maddelik bir kontrol listesi, Kapoor ve Narayanan~[17]
raporlama formları sunar; hiçbiri belirli bir operatöre uygulanıp sızdırıp
sızdırmadığına karar veremez. Bölüm~5.4 uygulanabilen ve hüküm yerine profil
raporlayan bir sertifika tanımlar.

\emph{Üçüncüsü, büyüklük yalnızca ham hata farkı olarak raporlanmıştır.} Ham fark
nedensel hataya bölünür; nedensel boru hattı zayıfken büyür ve farklı tavanlara
sahip operatörler kıyaslanamaz. Bölüm~5.3 farkı ulaşılabilir, operatöre göreli
bir oracle'a çapalar.

Kurulum ayrıca bu dergi açısından önemli bir biçimde farklıdır: yukarıdaki tüm
ölçümler saatlik çözünürlükte toplam ya da ulusal serilerdedir; oysa yedek ve
kısıtlama kararları 10--30 dakikalık türbin düzeyi SCADA üzerinde alınır. Türbin
düzeyi gücün uzamsal düzleştirmesi yoktur, oynaklığı çok daha yüksektir, pitch ve
kısıtlama etkileri taşır. Bu nedenle ölçümü türbin düzeyinde, tek bir nedensel
protokol altında, türbinler, ufuklar, rastgele başlangıçlar ve tahminci aileleri
boyunca raporluyoruz.

\subsection{Katkılar}

\begin{enumerate}
\item Bantları yalnız eğitim döneminden öğrenilen, kayan VMD'den istatistiksel
olarak ayırt edilemeyen nedensel doğruluğa yaklaşık 1300 kat düşük güncelleme
gecikmesiyle ulaşan \textbf{kilitli nedensel spektral filtre bankası (LCSF)}
(Bölüm~4.1 ve~6.5). Daha doğru bir tahminci olarak sunulmamaktadır; değeri
dürüst protokolü uygulanabilir kılmasıdır.
\item \textbf{Eğitim gerektirmeyen nedensellik sertifikası}: gelecek-bozma
duyarlılık profili $D(\tau;\delta)$ ve önek değişmezliği testi; sızıntı
\emph{erişimi} ve \emph{kazancı} raporlar. Sızdıran operatörleri birbirinden
ayırır --- tam seri VMD'nin duyarlılığı gecikmeyle azalmazken iki yönlü
süzmeninki azalır (Bölüm~5.4 ve~6.3).
\item Ulaşılabilir, operatöre göreli bir oracle'a çapalanmış
\textbf{oracle-normalize sızıntı oranı} $\lambda$; iki sızdıran operatörü ham
farkın tersine sıralar (Bölüm~5.3 ve~6.4).
\item Altı türbin, altı ufuk, sekiz tohum ve dört tahminci ailesi boyunca
Diebold--Mariano testiyle \textbf{türbin düzeyi ölçüm}, yerel 10 dakikalık
ızgara dahil.
\item \textbf{Olumsuz sonuçların tam raporlanması}: nedensel protokolde
ayrıştırma ham SCADA girdilerine göre tutarlı kazanç sağlamaz ve nedensel
modellerin persistence karşısındaki konumu eğitim kaybına bağlıdır.
\end{enumerate}

% =====================================================================
\section{İlgili çalışmalar}
% =====================================================================
\tasi{İngilizce sürüm Bölüm 2 ile aynı yapı. En yakın öncül iş (Jeon ve ark.,
Zhang ve ark., Zhou ve ark., Valdivieso ve ark.) başa alınacak; ardından
sızıntı-önleyici tasarımlar, rüzgâr gücü değerlendirme pratiği ve nedensel/
çevrimiçi ayrıştırma. Konumlandırma tablosunun son sütunu
\emph{sunulan operasyonel araç} olacak --- ayrımın yapıldığı yer burası.}

% =====================================================================
\section{Veri ve protokol}
% =====================================================================
\tasi{Taslak v2 §3.1--3.4 tamamı. SDWPF, 1500 kW anma gücü, örnekleme aralığı
bulgusu ve 30 dk ortak ızgara; pencereleme, ufuklar; 37 öznitelik; kronolojik
\%70/15/15 bölme, 12 adım purge, yalnız eğitimden ölçekleme.}
\hesapla{rolling-origin değerlendirme ya da en azından aylık alt-dönem sonuçları}
\hesapla{ikinci çiftlik: Kelmarsh (6 türbin, 10 dk SCADA)}

% =====================================================================
\section{Ayrıştırma operatörleri}
% =====================================================================
\tasi{Taslak v2 §4.1--4.2 ve §4.4: tam seri VMD (sızdıran), kayan VMD
(nedensel) ve tahminci tanımları.}

\subsection{Kilitli nedensel spektral filtre bankası}
\tasi{Taslak v2 §4.3.} VMD yalnız eğitim kesitine uygulanarak $K$ merkez
frekansı elde edilir; ardışık merkez frekansların geometrik ortalamaları bant
sınırı olur ve üçüncü dereceden Butterworth filtrelerinden oluşan bir banka
(bir alçak geçiren, iki bant geçiren, bir yüksek geçiren) kesin olarak ileri
yönde uygulanır; başlangıç durumu ilk gözlemden türetilir. Aynı bantların sıfır
fazlı uygulanması, nedensel olandan yalnızca süzme yönüyle ayrılan ve bu nedenle
kontrol operatörü işlevi gören sızdıran karşılığı verir.

\textbf{Gerekli eklemeler.} (i) Yapıyı VMD'nin filtre bankası okumasına~[2]
dayandır ve önsel merkez frekanslarıyla genelleştirilmiş VMD'den~[18]
farklandır. (ii) LCSF'nin test anında VMD'nin matematiksel eşdeğeri olarak
sunulmadığını belirt. (iii) İkinci dereceden bölüm kontrolünü raporla (taslak v2
§6.17): iki uygulama arasındaki en büyük fark bileşen başına en çok
$8\times10^{-6}$ kW, yani filtrelerin sayısal biçimi sonuçları etkilemiyor.

% =====================================================================
\section{Değerlendirme çerçevesi}
% =====================================================================
% [YENİ] — makalenin metodolojik çekirdeği

\subsection{Hata ölçütleri ve anlamlılık}
\tasi{Taslak v2 §5: MAE, RMSE, nMAE, nRMSE, $R^2$; referansa göre skill;
Bartlett uzun dönem varyansı ve Harvey--Leybourne--Newbold küçük örneklem
düzeltmesiyle Diebold--Mariano; Holm düzeltmesi; blok bootstrap.}

\subsection{Referans modeller}
\textsf{[YENİ]} Yalnız persistence topluluk referansı değildir. Nielsen ve
ark.~[19] MSE-optimal birleşimini benimsiyoruz:
\begin{equation}
\hat p(t+k) = a_k\,p_t + (1-a_k)\,\bar p,
\qquad
a_k = \frac{\frac{1}{N}\sum_t \tilde p_t \tilde p_{t+k}}
            {\frac{1}{N}\sum_t \tilde p_t^{\,2}},
\qquad \tilde p_t = p_t-\bar p,
\label{eq:nielsen}
\end{equation}
burada $k$ gecikmeli otokorelasyon $a_k$ ve ortalama $\bar p$ yalnız eğitim
döneminden ve o yazarların önerdiği gibi türbin bazında kestirilir.
$a_k\to1$'de persistence'a, $a_k\to0$'da climatology'ye indirgenir. Birleşim
kare hatayı minimize ettiği için MAE'si persistence'ınkini aşabilir; bu nedenle
iki ölçüt de raporlanır.

\subsection{Sızıntı oranları}
Bir operatörün eşleştirilmiş nedensel ve sızdıran sürümleri için ham fark
\begin{equation}
\mathrm{LSI} =
\frac{E_{\mathrm{nedensel}}-E_{\mathrm{sızdıran}}}{E_{\mathrm{nedensel}}}\times 100
\end{equation}
nedensel hataya böler; nedensel boru hattı zayıfken --- ki bir sızıntı
çalışmasının bulunması muhtemel rejim budur --- şişer ve farklı tavanlı
operatörler arasında kıyaslanamaz.
\textsf{[YENİ]} Bu nedenle oracle-normalize bir oran da raporluyoruz.
$E_{\mathrm{oracle}}$, \emph{aynı} operatörün hedefi içeren bir pencerede
uygulanmış hâlinin hatası olsun; yani bu operatörün cevabı görmesine izin
verilseydi yapabileceğinin en iyisi. O zaman
\begin{equation}
\lambda =
\frac{E_{\mathrm{nedensel}}-E_{\mathrm{sızdıran}}}
     {E_{\mathrm{nedensel}}-E_{\mathrm{oracle}}}
\label{eq:lambda}
\end{equation}
şunu yanıtlar: bu operatörün kabul ettiği sızıntının ne kadarını denetlenen
protokol gerçekleştiriyor? Oracle çapalama tahmin doğrulamasında yerleşiktir~[20];
mükemmel çapanın çoğu zaman ulaşılamaz olduğu uyarısıyla birlikte.
Denk.~\eqref{eq:lambda} bunun yerine ulaşılabilir, operatöre göreli bir oracle
kullanır. Her raporlanan değerle birlikte iki sınır belirtilmelidir:
$E_{\mathrm{nedensel}}\le E_{\mathrm{oracle}}$ ise $\lambda$ tanımsızdır; ve
kayan VMD'nin ileri-bakış eğrisi monoton olmadığından $E_{\mathrm{oracle}}$
kesin bir üst sınır değil, tanımlanmış bir referans noktasıdır. $\lambda$ ham
farkı tamamlar, yerine geçmez.

\subsection{Nedensellik sertifikası}
\textsf{[YENİ]} $\Phi$, $x$ serisini $K$ bileşene eşlesin;
$\Phi(x)_{t,k}$ $t$ anındaki $k$. bileşen ve
$\mathcal F_t=\sigma(x_s: s\le t)$ olsun. Operatör, her $t$ ve $k$ için
$\Phi(x)_{t,k}$ $\mathcal F_t$-ölçülebilir ise \emph{nedenseldir}.

Nedensellik hesaplamanın bir özelliğidir ve verinin erişilebilirliği veri
\emph{değerlerine} bağlı olmayan boru hatları için statik analizle karara
bağlanabilir. Tam seri VMD, kayan VMD ve EMD ailesi bu türden değildir: her biri
$t$ anındaki çıktısı tüm girdi penceresinin değerlerine bağlı olan yinelemeli
bir optimizasyondur. Bu sınıf için deneysel ve operatörden bağımsız bir prob
elimizdeki araçtır ve neyi ortaya koyduğunu belirtiyoruz: denenen küme üzerinde
gerekli bir koşul, asla sızıntının yokluğu değil.

\paragraph{Test 1: gelecek-bozma profili.}
Bozma anı $t_0$, genlik $\delta$ ve gecikme $\tau\ge0$ için
\begin{equation}
D_\Phi(\tau;\delta)=
\max_{t\le t_0-\tau}\;\max_k\;
\bigl|\Phi(x)_{t,k}-\Phi\!\left(x\oplus\delta\mathbf 1_{\{s>t_0\}}\right)_{t,k}\bigr| .
\end{equation}
$\Phi$ nedensel ise $D_\Phi\equiv0$'dır. Profilin raporlanması, $D_\Phi>0$ olan
en büyük denenen $\tau$ olan \emph{sızıntı erişimini} ve $D_\Phi$'nin $\delta$
ile büyümesi olan \emph{sızıntı kazancını} verir.

\paragraph{Test 2: önek değişmezliği.}
Bağımsız olarak, rastgele kesim noktalarında
$\Phi(x_{1:t})_{t,k}=\Phi(x_{1:N})_{t,k}$ aranır. Test~1 geleceğin bozulmasına
duyarlılığı, Test~2 geleceğin varlığına bağımlılığı yoklar.

\paragraph{Ortaya koymadıkları.} Geçmek nedenselliği kanıtlamaz; denenen anlar,
gecikmeler ve genliklerde bağımlılık saptanmadığını ortaya koyar ve bu
parametrelerle birlikte raporlanmalıdır.

\subsection{Ağırlıklı sızıntı maruziyeti ve dereceli ölçüt}
% [YENİ — HESAPLA]

Denk.~(5) bileşenler üzerinde maksimum alır; yani yalnızca operatörün sızdırıp
sızdırmadığını yanıtlar. Pratikte iki soru daha önemlidir: sızıntıyı
\emph{hangi} bileşenler taşıyor ve tahminci onları fiilen kullanıyor mu?
Birincisi çalışılmıştır: Zhang ve ark.~[15] ampirik mod fonksiyonlarını
tahminciye teker teker ekleyip yapaylığın ağırlıkla yüksek frekanslı
bileşenlerce taşındığını raporlar. Ancak o ölçü iki ayrı niceliği birleştirir ---
operatörün bir bileşende ne kadar gelecek açığa çıkardığı ile modelin bunu ne
kadar kullandığı --- çünkü tahmin hatasındaki değişimden okunur. Sertifika
bunları ayırır, zira $D_k$ yalnızca $\Phi$'nin bir özelliğidir.

Denk.~(5)'i bileşen bazında çözelim,
\begin{equation}
D_{\Phi,k}(\tau;\delta)=
\max_{t\le t_0-\tau}
\bigl|\Phi(x)_{t,k}-\Phi\!\left(x\oplus\delta\mathbf 1_{\{s>t_0\}}\right)_{t,k}\bigr| ,
\label{eq:profilek}
\end{equation}
ve $w_k$ tahmincinin $k$. bileşene dayanma payı,
$\hat w_k = w_k/\sum_j w_j$ normalize hâli olsun (Ek C'de zaten hesaplanmış üç
ölçütü de raporluyoruz: öngörü-değeri-değişimi, permütasyon $\Delta$MAE ve
ortalama mutlak SHAP). Her bileşeni kendi standart sapması $\sigma_k$ ile
ölçekleyerek toplamı boyutsuzlaştırırsak, \emph{ağırlıklı sızıntı maruziyeti}
\begin{equation}
\Lambda(\tau;\delta)=\sum_{k=1}^{K}\hat w_k\,\frac{D_{\Phi,k}(\tau;\delta)}{\sigma_k}.
\label{eq:exposure}
\end{equation}
$\Lambda=0$ ancak ve ancak operatör nedenselse ya da model sızdıran her bileşeni
görmezden geliyorsa sağlanır. $\{\hat w_k D_{\Phi,k}/\sigma_k\}$ kümesinin
yoğunlaşması, sızıntının tek bir bileşende mi toplandığını yoksa bankaya mı
dağıldığını gösterir.

\paragraph{Dereceli kabul ölçütü.} Yukarıdaki hâliyle sertifika ikilidir:
operatör ancak $D$ tam olarak sıfırsa geçer. Bu kırılgandır; ihmal edilebilir
düzeyde sızdıran ve hiç kullanılmayan bir bileşen, yapaylığın tamamını taşıyan
bir bileşenle aynı muameleyi görür. Bunun yerine
\begin{equation}
\Lambda(0;\delta_{\mathrm{ref}})\le c
\label{eq:criterion}
\end{equation}
koşulunu öneriyoruz ve $c$'yi uygulayıcının zaten ölçebildiği bir niceliğe
kalibre ediyoruz: $c$, ima ettiği hata şişmesi aynı tahmincinin rastgele
başlangıçlar arası standart sapmasının altına düşen maruziyet düzeyidir (bu
veride 2,4--6,6 kW; sekiz tohum tablosu ve taslak v2 §6.3). Ölçüt şöyle okunur:
\emph{bir boru hattı, hâlâ kabul ettiği sızıntı hatayı yeniden eğitmekten daha
az oynattığında kabul edilebilir.} Hem $\Lambda$ hem ham profil raporlanır; okur
istediğini uygulayabilir.

\paragraph{İddianın kapsamı.} $w_k$ modele bağlıdır, tamsayı değildir ve yeniden
eğitimler arasında sabit değildir; dolayısıyla Denk.~\eqref{eq:criterion} bir
ağırlıklı \emph{kabul ölçütüdür}, ağırlıklı $k$-out-of-$n$ güvenilirlik modeli
değildir ve böyle sunulmaktadır. $\Lambda$ simgesi, Bölüm~5.3'teki
oracle-normalize oran $\lambda$'dan ayrıdır.

\hesapla{(i) \texttt{leakage/certificate.py::pertürbasyon\_profili} bileşen
bazında çözülecek --- şu an $k$ üzerinde indirgeme yapıyor --- ve betik 29
yeniden koşulacak; maliyeti tam seri VMD belirliyor, varsayılan prob ızgarasıyla
yaklaşık 1--2 saat. (ii) $\hat w_k$ mevcut Ek C önem tablosundan alınacak;
yeniden eğitim gerekmiyor. (iii) $c$, hâlihazırda raporlanan sekiz tohumlu MAE
standart sapmalarına kalibre edilecek.}

% =====================================================================
\section{Bulgular}
% =====================================================================
% Tüm tablolar taslak v2'den birebir aktarıldı.

\subsection{Sızıntıyla ilişkili farkın büyüklüğü}
\tasi{Taslak v2 §6.1--6.2 metni.}

\begin{table}[htbp]\centering
\caption{Türbin 57, çok ufuklu MAE (kW) ve LSI (tohum 42).}
\footnotesize
\begin{tabular}{lrrrrrr}
\toprule
Ufuk (sa) & Persistence & Ham LSTM & CatBoost Direct &
Kayan VMD (ned.) & Tam seri VMD (sızd.) & LSI-VMD (\%) \\
\midrule
0,5 &  77,80 &  90,04 &  78,78 &  86,76 & 53,29 & 38,6 \\
1   & 108,42 & 114,55 & 106,12 & 116,11 & 64,71 & 44,3 \\
2   & 153,78 & 164,28 & 143,35 & 158,98 & 67,73 & 57,4 \\
3   & 187,89 & 211,91 & 168,05 & 191,00 & 73,43 & 61,6 \\
4   & 214,98 & 212,26 & 187,36 & 215,86 & 96,81 & 55,2 \\
6   & 255,53 & 240,69 & 215,69 & 247,92 & 88,55 & 64,3 \\
\bottomrule
\end{tabular}
\end{table}

\begin{table}[htbp]\centering
\caption{Bir saatlik ufuk, altı türbin, tohum 42. MAE kW; parantez içinde
persistence'a göre Skill (\%). DM negatifse sızdıran model daha düşük hatalı;
tüm $p<10^{-15}$.}
\scriptsize
\begin{tabular}{lrrrrrrrrrr}
\toprule
& \multicolumn{5}{c}{Nedensel} & \multicolumn{2}{c}{Sızdıran} &
\multicolumn{2}{c}{LSI (\%)} & \\
\cmidrule(lr){2-6}\cmidrule(lr){7-8}\cmidrule(lr){9-10}
Türbin & Pers. & Ham LSTM & CatBoost & Kayan VMD & LCSF &
Tam VMD & LCSF iki y. & VMD & LCSF & DM (VMD) \\
\midrule
57  & 108,42 & 114,55 (-5,7)  & 106,12 (2,1)  & 116,11 (-7,1)  & 116,29 (-7,3)  & 64,71 (40,3) & 57,34 (47,1) & 44,3 & 50,7 & $-26,56$ \\
10  & 110,58 & 131,73 (-19,1) & 116,72 (-5,6) & 128,17 (-15,9) & 127,71 (-15,5) & 69,58 (37,1) & 55,57 (49,7) & 45,7 & 56,5 & $-27,82$ \\
83  &  95,85 & 114,43 (-19,4) &  98,84 (-3,1) & 111,64 (-16,5) & 111,34 (-16,2) & 61,67 (35,7) & 57,89 (39,6) & 44,8 & 48,0 & $-24,77$ \\
120 &  95,08 &  97,32 (-2,4)  &  93,37 (1,8)  &  95,59 (-0,5)  &  96,83 (-1,8)  & 63,42 (33,3) & 57,29 (39,7) & 33,7 & 40,8 & $-18,49$ \\
61  & 107,31 & 117,78 (-9,8)  & 107,11 (0,2)  & 119,48 (-11,3) & 119,26 (-11,1) & 70,15 (34,6) & 56,01 (47,8) & 41,3 & 53,0 & $-26,12$ \\
84  & 103,19 & 109,12 (-5,7)  & 102,40 (0,8)  & 116,34 (-12,7) & 110,15 (-6,7)  & 72,61 (29,6) & 55,03 (46,7) & 37,6 & 50,0 & $-24,66$ \\
\bottomrule
\end{tabular}
\end{table}

\begin{table}[htbp]\centering
\caption{LSI (\%), sekiz tohum (42--49) ortalama $\pm$ standart sapma.
Türbinler tek çiftliktedir ve bağımsız tekrar sayılmaz.}
\footnotesize
\begin{tabular}{lrr}
\toprule
Türbin & LSI-VMD & LSI-LCSF \\
\midrule
57  & $39,9\pm4,4$ & $48,7\pm5,0$ \\
10  & $45,1\pm3,1$ & $53,1\pm3,8$ \\
83  & $45,2\pm2,9$ & $52,0\pm4,5$ \\
120 & $37,7\pm5,2$ & $47,0\pm5,2$ \\
61  & $42,8\pm5,7$ & $50,4\pm4,9$ \\
84  & $40,7\pm4,7$ & $50,3\pm2,9$ \\
\midrule
Türbinler arası ort. & $41,9\pm3,0$ & $50,3\pm2,2$ \\
\bottomrule
\end{tabular}
\end{table}

\subsection{Mekanizma: ileri-bakış doz--tepkisi}
\tasi{Taslak v2 §6.10--6.11, sıkıştırılmış.}

\begin{table}[htbp]\centering
\caption{Türbin 57, 1 sa, üç tohum: bileşene tanınan ileri-bakış miktarı $m$'ye
göre MAE (tohum ortalaması) ve LSI-LCSF (ortalama $\pm$ std). $m\ge H=2$
koşulları hedefi içerir; gerçekçi kullanım değil, üst sınır stres testidir.}
\footnotesize
\begin{tabular}{lrrr}
\toprule
İleri-bakış $m$ (adım) & Gelecek (dk) & MAE (kW) & LSI-LCSF (\%) \\
\midrule
Nedensel LCSF (tek yönlü, referans) & 0 & 115,94 & --- \\
0 (sıfır fazlı, gelecek yok)        & 0 & 113,27 & $2,3\pm1,4$ \\
1                                   & 30 & 85,86 & $26,0\pm1,8$ \\
2 ($=H$, hedef bileşende)           & 60 & 11,05 & $90,5\pm1,0$ \\
4                                   & 120 & 29,10 & $74,9\pm0,6$ \\
8                                   & 240 & 59,43 & $48,7\pm6,3$ \\
16                                  & 480 & 63,50 & $45,2\pm2,7$ \\
48                                  & 1440 & 61,60 & $46,8\pm5,1$ \\
$\infty$ (tam seri, iki yönlü)      & --- & 61,73 & $46,7\pm7,7$ \\
\bottomrule
\end{tabular}
\end{table}

Sıfır fazlı ama geleceğe bakmayan süzme farkı yalnızca \%2,3'e çıkarıyor; yani
iki yönlü ile tek yönlü süzme arasındaki fark genlik veya faz yapaylığı değil.
Bir adım ileri-bakış \%26,0'a, hedefin dahil edilmesi \%90,5'e çıkarıyor. Hiç
gelecek tanımadan bileşeni pencere sonundan uzaklaştıran geriye kaydırılmış
nedensel kontrol MAE'yi yaklaşık \%1 değiştiriyor (taslak v2 §6.11), yani sınır
etkisi bu veride küçük.

\subsection{Nedensellik sertifikası}
\tasi{Taslak v2 §6.18 metni.} Prob kümesi: $\tau\in\{0,1,2,4,8,16,32\}$ adım,
$\delta\in\{50,100,200,400\}$ kW, 20 bozma anı (tam seri VMD için 5) ve
50 rastgele kesim noktası (tam seri VMD için 10).

\begin{table}[htbp]\centering
\caption{Nedensellik sertifikası, Türbin 57. $D(\tau;\delta=200$ kW$)$ kW
cinsinden. Erişim, $D>0$ olan en büyük denenen $\tau$; 32 adım ızgara sınırıdır.}
\footnotesize
\begin{tabular}{lrrrlrl}
\toprule
Operatör & $D(\tau{=}0)$ & $D(\tau{=}8)$ & $D(\tau{=}32)$ & Erişim &
Önek maks. fark (kW) & Önek ihlali \\
\midrule
LCSF, nedensel            &   0,00 &   0,00 &  0,00 & yok       & 0 & 0 / 50 \\
Kayan VMD, nedensel       &   0,00 &   0,00 &  0,00 & yok       & 0 & 0 / 50 \\
LCSF, iki yönlü (sızd.)   &  87,21 &  55,39 & 10,17 & $\ge32$   & 1.060 & 50 / 50 \\
Tam seri VMD (sızd.)      & 113,49 & 235,44 & 96,57 & $\ge32$   & 386 & 10 / 10 \\
\bottomrule
\end{tabular}
\end{table}

Her iki nedensel operatör de denenen her $\tau$ ve $\delta$'da tam olarak sıfır
veriyor ve önek ihlali göstermiyor. İki sızdıran operatör niteliksel olarak
ayrışıyor: iki yönlü süzmenin duyarlılığı, sonlu destekli bir operatörden
beklendiği gibi gecikmeyle azalıyor; tam seri VMD'ninki azalmıyor. Bu, serinin
herhangi bir yerindeki değişiklikte tüm ayrıştırmayı yeniden dağıtan küresel bir
varyasyonel optimizasyonla tutarlıdır. Profil böylece yalnızca sızıntının
varlığını değil \emph{türünü} de ayırt ediyor.

\subsection{Oracle-normalize oran}
\tasi{Taslak v2 §6.14 metni, çekinceleriyle birlikte.}

\begin{table}[htbp]\centering
\caption{Türbin 57, 1 sa, üç tohum: hata ölçeğinin dört çapası ve $\lambda$.
$E_{\mathrm{oracle}}$, $m=H=2$ hedef-içeren koşuldur. Hatalar kW.}
\footnotesize
\begin{tabular}{lrrrrr}
\toprule
Operatör & $E_{\mathrm{nedensel}}$ & $E_{\mathrm{sızdıran}}$ &
$E_{\mathrm{oracle}}$ & LSI (\%) & $\lambda$ \\
\midrule
İki yönlü LCSF  & 115,94 & 61,73 & 11,05 & 46,7 & 0,52 \\
Tam seri VMD    & 116,78 & 68,66 & 64,64 & 41,1 & 0,92 \\
\bottomrule
\end{tabular}
\end{table}

İki ölçüt operatörleri ters sıralıyor: ham fark iki yönlü LCSF'de daha büyükken,
$\lambda$ tam seri VMD'nin operatörünün kabul ettiği sızıntının neredeyse
tamamını, iki yönlü süzmenin ise yaklaşık yarısını gerçekleştirdiğini gösteriyor.
Ham fark operatörleri farklı tavanlara karşı kıyaslıyor. Bu tek türbin, tek ufuk
ve üç tohumdur; tam seri VMD'nin oracle'ı zayıftır (68,66'ya karşı 64,64 kW),
dolayısıyla $\lambda$ bu hücrede duyarlıdır.
\hesapla{kalan türbin ve ufuklar için $\lambda$}

\subsection{Nedensel protokolün maliyeti}
\tasi{Taslak v2 §6.4 metni.}

\begin{table}[htbp]\centering
\caption{Tahmin anı gecikmesi, Türbin 57, tek çekirdek; kayan VMD için seri
boyunca 200 konumun ortancası.}
\footnotesize
\begin{tabular}{lrr}
\toprule
Yöntem & Tek tahmin anı & Test kümesi (5.252 nokta) \\
\midrule
Kayan VMD ($W=672$)        & 89,9 ms & $\approx$ 8 dk \\
LCSF (çevrimiçi güncelleme) & 69 µs  & 0,36 sn \\
\bottomrule
\end{tabular}
\end{table}

Nedensel doğruluk 1 ve 6 saatte istatistiksel olarak ayırt edilemiyor, 3 saatte
küçük ama anlamlı bir LCSF üstünlüğü var. Oran uygulamaya ve donanıma bağlıdır;
sonuç büyüklük mertebesidir ($10^3$). Çiftlik ölçeğinde sonucu doğrudandır: kayan
VMD tüm seri için türbin başına saat mertebesinde süre ister; 134 türbin için bu,
tahmin anı başına yeniden ayrıştırmayı uygulanamaz kılar. Kilitli banka bunu
saniyeye indirir.

\subsection{Nedensel modellerin referans modellere karşı durumu}
\tasi{Taslak v2 §6.15 metni; Nielsen referansının MAE'de daha zayıf referans
olduğu ve ona göre pozitif skill'in persistence'a göre olandan güçlü
okunmaması gerektiği çekincesiyle.}

\begin{table}[htbp]\centering
\caption{Bir saatlik ufuk, altı türbin: referans hataları (kW) ve nedensel
modellerin Nielsen referansına (Denk.~\eqref{eq:nielsen}) göre MAE skill'i (\%).}
\scriptsize
\begin{tabular}{lrrrrrrl}
\toprule
Türbin & $a_k$ & MAE pers. & MAE Nielsen & MAE clim. & RMSE pers. &
RMSE Nielsen & Nielsen'e göre Skill: LSTM / VMD / LCSF / CatBoost \\
\midrule
T57  & 0,837 & 108,41 & 125,67 & 338,07 & 190,21 & 186,37 & 8,8 / 7,6 / 7,5 / 15,6 \\
T10  & 0,847 & 110,58 & 135,29 & 405,95 & 199,52 & 198,35 & 2,6 / 5,3 / 5,6 / 13,7 \\
T83  & 0,837 &  95,85 & 122,54 & 371,01 & 181,64 & 182,63 & 6,6 / 8,9 / 9,1 / 19,3 \\
T120 & 0,819 &  95,08 & 112,11 & 292,29 & 169,21 & 165,78 & 13,2 / 14,7 / 13,6 / 16,7 \\
T61  & 0,826 & 107,31 & 127,02 & 335,40 & 190,36 & 186,82 & 7,3 / 5,9 / 6,1 / 15,7 \\
T84  & 0,833 & 103,19 & 119,89 & 322,52 & 180,62 & 177,20 & 9,0 / 3,0 / 8,1 / 14,6 \\
\bottomrule
\end{tabular}
\end{table}

\subsection{Eğitim kaybına duyarlılık}
\tasi{Taslak v2 §6.16 metni.}

\begin{table}[htbp]\centering
\caption{Türbin 57, 1 sa, mutlak hata kaybıyla eğitilmiş LSTM, üç tohum
(42--44). Persistence MAE'si 108,41 kW.}
\footnotesize
\begin{tabular}{llrr}
\toprule
Model & Protokol & MAE (kW) & Persistence'a göre MAE skill (\%) \\
\midrule
Ham LSTM              & nedensel             & $106,65\pm0,33$ & 1,6 \\
Kayan VMD--LSTM       & nedensel             & $107,22\pm0,31$ & 1,1 \\
LCSF--LSTM            & nedensel (tek yönlü) & $107,19\pm0,76$ & 1,1 \\
Tam seri VMD--LSTM    & sızdıran             & $61,61\pm3,10$  & --- \\
LCSF--LSTM            & sızdıran (iki yönlü) & $56,86\pm3,49$  & --- \\
\bottomrule
\end{tabular}
\end{table}

Eşleştirilmiş kayıpla nedensel modeller 106,65--107,22 kW'a ulaşıyor ve
persistence'ı \%1,1--1,6 geçiyor; sızıntıyla ilişkili fark değişmiyor
(LSI-VMD $42,5\pm2,7$\%, LSI-LCSF $46,9\pm3,4$\%). İki sonuç çıkıyor: ana bulgu
kare hata kaybının yapaylığı değildir, buna karşılık nedensel modellerin MAE'de
persistence karşısındaki konumu kaybın yapaylığıdır. Ayrıştırma kazancının
yokluğu mutlak hata kaybında da geçerli --- üç nedensel model birbirinden
yaklaşık 0,6 kW içinde kalıyor.

\subsection{Sağlamlık}
\tasi{Taslak v2 §6.3, §6.7--6.9, §6.12--6.13 ve §7.1'deki eşit eğitim penceresi
kontrolü; ayrıntı Ek B'de.}

% =====================================================================
\section{Tartışma}
% =====================================================================
\tasi{Taslak v2 §7, ham fark yerine maliyet ekseni etrafında yeniden
yazılacak.}

\subsection{Sınırlılıklar}
\tasi{Taslak v2 §7.1 neredeyse olduğu gibi. Eklenecek: sertifika denenen küme
üzerinde gerekli bir koşuldur; $\lambda$ oracle tanımına bağlıdır ve VMD
oracle'ı zayıftır; literatür taraması özet düzeydedir.}

% =====================================================================
\section{Sonuç}
% =====================================================================
\tasi{Taslak v2 §8, üç araç ve maliyet sonucu etrafında yeniden yazılacak.}

% =====================================================================
\section*{Beyanlar}
\tasi{Taslak v2. CRediT, çıkar çatışması, üretken yapay zekâ beyanı ve veri
erişilebilirliği. Yazar sırası tutarsızlığı çözülecek.}

\textbf{Veri erişilebilirliği --- gönderimden önce:} depo şu an özeldir ve
makaledeki adres okur için çözülmemektedir; herkese açık yapılmalı, kayıtlı
sonuç tabloları eklenmeli (depo şu an yalnız kod taşıyor) ve kalıcı bir Zenodo
DOI'si oluşturulup burada atıf verilmelidir.

% =====================================================================
\section*{Kaynakça}
\textsf{Kaynakça tek listeden, İngilizce sürümde yönetilmektedir
(\texttt{manuscript/en/RE\_manuscript\_v2.tex}). Yalnızca hakemli kaynaklar ana
listeye girer; preprint'ler işaretlidir.}

\end{document}
